% Original scattering bridge: fully solvable Lippmann--Schwinger check.
\subsection*{Worked example: outgoing Green function, Born series, and exact scattering}
The formal integral equation becomes much more concrete in
one dimension. Work in units $\hbar^2/(2m)=1$, choose
positive energy $E=k^2$ with $k>0$, and consider
\begin{equation}
\left(-\frac{\dd^2}{\dd x^2}+g\delta(x)\right)
\psi(x)=k^2\psi(x).
\end{equation}
The incident wave comes from the left. Outside $x=0$
the stationary equation is free, so the most general
outgoing scattering ansatz is
\begin{equation}
\psi(x)=
\begin{cases}
e^{\ii kx}+r e^{-\ii kx},&x<0,\\
t e^{\ii kx},&x>0.
\end{cases}
\label{eq:worked-delta-ansatz}
\end{equation}
Continuity of the wavefunction gives $\psi(0^-)=\psi(0^+)$,
hence $1+r=t$. Next integrate the differential equation
from $-\epsilon$ to $+\epsilon$:
\begin{equation}
-\psi'(\epsilon)+\psi'(-\epsilon)+g\psi(0)
=k^2\int_{-\epsilon}^\epsilon\psi(x)\dd x.
\end{equation}
The right side tends to zero as $\epsilon\to0$;
the delta function contributes $g\psi(0)$. Therefore
\begin{equation}
\psi'(0^+)-\psi'(0^-)=g\psi(0).
\end{equation}
The two derivatives obtained from
Eq.~\eqref{eq:worked-delta-ansatz} are
$\psi'(0^+)=\ii kt$ and $\psi'(0^-)=\ii k(1-r)$.
Use $t=1+r$ to find
\begin{equation}
\ii k[(1+r)-(1-r)]=g(1+r),
\qquad 2\ii kr=g(1+r).
\end{equation}
Solving the last equation step by step,
\begin{equation}
(2\ii k-g)r=g,\qquad
r=\frac{g}{2\ii k-g}
=\boxed{\frac{-\ii u}{1+\ii u}},
\quad
t=1+r=\boxed{\frac1{1+\ii u}},
\quad u=\frac{g}{2k}.
\label{eq:worked-delta-amplitudes}
\end{equation}
The outgoing current carried by each plane wave is proportional
to its signed wavenumber. Incoming and outgoing wavenumber
magnitudes are identical, so reflection and transmission
probabilities are simply squared amplitudes:
\begin{equation}
R=|r|^2=\frac{u^2}{1+u^2},\qquad
T=|t|^2=\frac1{1+u^2},\qquad
\boxed{R+T=1}.
\end{equation}
This current-conservation identity is an independent check.

\paragraph{Recover the same result from the integral equation.}
The one-dimensional outgoing free Green function for
$k^2+\partial_x^2$ is
\begin{equation}
G_0^+(x,x';k)=-\frac{\ii}{2k}e^{\ii k|x-x'|}.
\end{equation}
Its first derivative jumps by $1$ in the equation
$(k^2+\partial_x^2)G_0^+=\delta(x-x')$,
which can be checked on either side of $x=x'$.
The Lippmann--Schwinger equation becomes
\begin{equation}
\psi(x)=e^{\ii kx}+
g\,G_0^+(x,0;k)\psi(0).
\end{equation}
Setting $x=0$ gives an \emph{algebraic} equation:
\begin{equation}
\psi(0)=1-\frac{\ii g}{2k}\psi(0)
\quad\Longrightarrow\quad
\psi(0)=\frac1{1+\ii u}.
\end{equation}
For $x<0$, the scattered term is
$-\ii u\,\psi(0)e^{-\ii kx}$;
therefore $r=-\ii u/(1+\ii u)$,
exactly as in Eq.~\eqref{eq:worked-delta-amplitudes}.
For $x>0$ the total amplitude is
$1-\ii u\psi(0)=1/(1+\ii u)=t$.

\paragraph{Show precisely where the Born series converges.}
At the origin, the repeated interaction gives a geometric series
\begin{equation}
\psi(0)=\frac1{1+\ii u}
=\sum_{n=0}^{\infty}(-\ii u)^n,
\qquad |u|<1.
\end{equation}
The order-$N$ reflection amplitude is
\begin{equation}
r_N=-\ii u\sum_{n=0}^{N-1}(-\ii u)^n.
\end{equation}
Subtract it from the exact geometric sum:
\begin{equation}
r-r_N=r(-\ii u)^N,\qquad
|r-r_N|=|r|\,|u|^N.
\end{equation}
For $|u|>1$, successive Born terms grow, while the exact
amplitude remains bounded and satisfies $R+T=1$.
The radius of convergence here is a property of this
specified delta-potential series; other potentials require
their own analysis. The complete
\href{https://github.com/KangqiaoLiu/advanced-quantum-mechanics-lecture-note/blob/main/companion/labs/IC-05-03/index.html}
{interactive scattering problem} illustrates both behaviors.
